\chapter{Continuous Models}

\section{Core assumption}
\begin{defbox}
A continuous model represents system evolution over a continuum of time, so state changes are described using trajectories, rates, and dynamics rather than only discrete steps.
\end{defbox}

\section{Typical properties}
\begin{itemize}
  \item State changes smoothly over time.
  \item Derivatives, rates, and invariants are often central.
  \item The model is common in control and physical systems.
\end{itemize}

\section{Why it matters}
\begin{keybox}
Continuous models are useful when the system interacts with physical processes, where discrete event abstraction is too coarse.
\end{keybox}

\section{Exam prompts}
\begin{itemize}
  \item Explain the difference between continuous and discrete-time reasoning.
  \item Give an example of a system naturally modeled continuously.
  \item Describe what it means for a property to hold over an interval of time.
\end{itemize}
